% =====================================================================
% Figure contract
% 核心结论：§10 的十个开放方向可归纳为四个互补研究集群，并由统一的证据标签约束解读。
% 图原型：mode C / 中等复杂度 2×2 卡片网格；无流程连线、无中心 hero。
% 证据层级：四张等尺寸集群卡为主；十个方向 chip 为证据；底部注释条限定术语含义。
% 能不能更少？：不能；删卡会丢失集群，删 chip 会丢失指定方向，删注释条会丢失证据边界。
%
% Step ① 设计文档（Form A）
% 1 领域：Computer-Use Agents 中文综述；技术术语保留英文。
% 2 格式：TikZ standalone，XeLaTeX + fontspec + xeCJK；全图 sans serif。
% 3 目标印刷宽度 W_print = 18.0 cm；目标印刷最小字号 = 7.5 pt。
%   设计画布 W_ink 约 18.0 cm，最小正文 7.6 pt，不依赖后期缩放放大字号。
% 4 布局：标题 + 2×2 等尺寸网格 + 底部整宽注释条；阅读方向为左上到右下。
% 5 模块：4 张 8.70×4.20 cm 卡片、10 个单行 chip、1 个整宽注释条。
% 6 连线：无；四个集群是并列分类，不暗示顺序、因果或依赖关系。
% 7 空间：列间距 0.60 cm、行间距 0.58 cm；卡内标题与 chip 均按锚点定位。
% 8 强调：集群名用粗体与深色边框；chip 用同色系更浅底色；注释条使用中性灰。
% 9 密度：中等；纯分类结构无数值，不嵌入 chart，避免为装饰引入未被正文支持的信息。
% 10 ASCII 草图：
%      ┌──────── 数据与验证 ────────┐  ┌──────── 长程与鲁棒 ────────┐
%      │ chip / chip / chip         │  │ chip / chip / chip         │
%      └────────────────────────────┘  └────────────────────────────┘
%      ┌──────── 交互与适应 ────────┐  ┌──────── 部署与人因 ────────┐
%      │ chip / chip                │  │ chip / chip                │
%      └────────────────────────────┘  └────────────────────────────┘
%      ┌────────────────── 证据术语注释条 ──────────────────────────┐
%      └─────────────────────────────────────────────────────────────┘
%
% Pre-flight：P1 设计文档已写；P2 四卡由 positioning 按构造严格对齐；
% P3 N/A（无连线）；P4 N/A（无嵌入数值图）；P5 N/A（无 rail）；
% P6 固定尺寸卡片与注释条；P7 已对照 lessons 的 CJK、字号与间距规则。
% =====================================================================

\documentclass[border=8pt]{standalone}
\usepackage{fontspec}
\usepackage{xeCJK}
\usepackage{xcolor}
\usepackage{tikz}
\setCJKmainfont{Heiti SC}
\setCJKsansfont{Heiti SC}
\usetikzlibrary{positioning,calc}

% Low-saturation palette shared with cua-survey-fig1-overview.tex.
\definecolor{ink}{HTML}{27364A}
\definecolor{muted}{HTML}{748092}
\definecolor{blueLine}{HTML}{6D879E}
\definecolor{blueFill}{HTML}{E7EEF4}
\definecolor{purpleLine}{HTML}{817594}
\definecolor{purpleFill}{HTML}{EEEAF2}
\definecolor{tealLine}{HTML}{648D8A}
\definecolor{tealFill}{HTML}{E5F0EE}
\definecolor{roseLine}{HTML}{9A777B}
\definecolor{roseFill}{HTML}{F3E9EA}
\definecolor{noteLine}{HTML}{A5ADB8}
\definecolor{noteFill}{HTML}{F2F4F6}

\tikzset{
  every node/.style={
    font=\sffamily\fontsize{8.8}{10.6}\selectfont,
    text=ink
  },
  cluster card/.style={
    rectangle,
    rounded corners=6pt,
    minimum width=8.70cm,
    minimum height=4.20cm,
    inner sep=0pt,
    outer sep=0pt,
    line width=0.75pt
  },
  cluster title/.style={
    anchor=west,
    align=left,
    inner sep=0pt,
    font=\sffamily\fontsize{10.7}{12.6}\selectfont\bfseries
  },
  chip/.style={
    rectangle,
    rounded corners=3pt,
    minimum width=7.66cm,
    minimum height=0.70cm,
    text width=7.18cm,
    align=left,
    inner xsep=0.24cm,
    inner ysep=0.13cm,
    outer sep=0pt,
    line width=0.42pt,
    font=\sffamily\fontsize{8.8}{10.4}\selectfont
  },
  note bar/.style={
    rectangle,
    rounded corners=4pt,
    draw=noteLine,
    fill=noteFill,
    minimum width=18.00cm,
    minimum height=1.62cm,
    text width=17.20cm,
    align=left,
    inner xsep=0.40cm,
    inner ysep=0.24cm,
    outer sep=0pt,
    line width=0.60pt,
    font=\sffamily\fontsize{7.6}{10.2}\selectfont
  }
}

\begin{document}
\sffamily
\begin{tikzpicture}

% Four fixed-size cards; center-to-center positioning makes both axes exact.
\node[cluster card,draw=blueLine,fill=blueFill] (data) at (0,0) {};
\node[cluster card,right=0.60cm of data,draw=purpleLine,fill=purpleFill] (robust) {};
\node[cluster card,below=0.58cm of data,draw=tealLine,fill=tealFill] (interact) {};
\node[cluster card,right=0.60cm of interact,draw=roseLine,fill=roseFill] (deploy) {};

% Title.
\coordinate (gridNorth) at ($(data.north)!0.5!(robust.north)$);
\node[
  anchor=south,
  text=ink,
  font=\sffamily\fontsize{13.6}{16.0}\selectfont\bfseries
] at ([yshift=0.48cm]gridNorth)
  {开放研究议程：四大集群十个方向（§10）};

% Card headers and understated separators.
\node[cluster title,text=blueLine] (dataTitle)
  at ([xshift=0.50cm,yshift=-0.50cm]data.north west) {数据与验证};
\draw[blueLine!55,line width=0.45pt]
  ([xshift=0.50cm,yshift=-0.96cm]data.north west) --
  ([xshift=-0.50cm,yshift=-0.96cm]data.north east);

\node[cluster title,text=purpleLine] (robustTitle)
  at ([xshift=0.50cm,yshift=-0.50cm]robust.north west) {长程与鲁棒};
\draw[purpleLine!55,line width=0.45pt]
  ([xshift=0.50cm,yshift=-0.96cm]robust.north west) --
  ([xshift=-0.50cm,yshift=-0.96cm]robust.north east);

\node[cluster title,text=tealLine] (interactTitle)
  at ([xshift=0.50cm,yshift=-0.50cm]interact.north west) {交互与适应};
\draw[tealLine!55,line width=0.45pt]
  ([xshift=0.50cm,yshift=-0.96cm]interact.north west) --
  ([xshift=-0.50cm,yshift=-0.96cm]interact.north east);

\node[cluster title,text=roseLine] (deployTitle)
  at ([xshift=0.50cm,yshift=-0.50cm]deploy.north west) {部署与人因};
\draw[roseLine!55,line width=0.45pt]
  ([xshift=0.50cm,yshift=-0.96cm]deploy.north west) --
  ([xshift=-0.50cm,yshift=-0.96cm]deploy.north east);

% Data and verification: three aligned chips.
\node[chip,anchor=north west,draw=blueLine!65,fill=white!68!blueFill]
  at ([xshift=0.52cm,yshift=-1.24cm]data.north west)
  {10.1 可扩展可验证数据};
\node[chip,anchor=north west,draw=blueLine!65,fill=white!68!blueFill]
  at ([xshift=0.52cm,yshift=-2.10cm]data.north west)
  {10.4 验证中心化 computer use};
\node[chip,anchor=north west,draw=blueLine!65,fill=white!68!blueFill]
  at ([xshift=0.52cm,yshift=-2.96cm]data.north west)
  {10.10 标准化与部署导向评测};

% Long-horizon and robustness: three aligned chips.
\node[chip,anchor=north west,draw=purpleLine!65,fill=white!68!purpleFill]
  at ([xshift=0.52cm,yshift=-1.24cm]robust.north west)
  {10.2 长程规划与状态追踪};
\node[chip,anchor=north west,draw=purpleLine!65,fill=white!68!purpleFill]
  at ([xshift=0.52cm,yshift=-2.10cm]robust.north west)
  {10.3 鲁棒 grounding 与动态环境};
\node[chip,anchor=north west,draw=purpleLine!65,fill=white!68!purpleFill]
  at ([xshift=0.52cm,yshift=-2.96cm]robust.north west)
  {10.5 错误检测与恢复};

% Interaction and adaptation: two chips centered in the same card body.
\node[chip,anchor=north west,draw=tealLine!65,fill=white!68!tealFill]
  at ([xshift=0.52cm,yshift=-1.52cm]interact.north west)
  {10.6 混合 GUI-CLI-API 交互};
\node[chip,anchor=north west,draw=tealLine!65,fill=white!68!tealFill]
  at ([xshift=0.52cm,yshift=-2.58cm]interact.north west)
  {10.7 个性化与持续学习};

% Deployment and human factors: two chips centered in the same card body.
\node[chip,anchor=north west,draw=roseLine!65,fill=white!68!roseFill]
  at ([xshift=0.52cm,yshift=-1.52cm]deploy.north west)
  {10.8 安全·隐私·人类监督};
\node[chip,anchor=north west,draw=roseLine!65,fill=white!68!roseFill]
  at ([xshift=0.52cm,yshift=-2.58cm]deploy.north west)
  {10.9 高效与端侧 CUA};

% Full-width evidence note.
\coordinate (gridBottomCenter) at ($(interact.south)!0.5!(deploy.south)$);
\node[note bar,anchor=north] (note) at ([yshift=-0.48cm]gridBottomCenter)
  {[Observed Tension] = 已有证据揭示的结构性冲突；[Validated Gap] = 已被
   benchmark/系统边界明确暴露但尚未闭合的能力缺口——后者不表示无人研究（§10.1 说明）};

\end{tikzpicture}
\end{document}
